\section{Binaurale 3D-Audio-Ausgabe mit HRTF}
\label{sec:binaural_3d_audio}

Mit der Erweiterung auf 3D-Raumgeometrie wurde die Audioausgabe um räumliche HRTF-basierte Verarbeitung erweitert.

\subsection{Richtungsdaten aus der Simulation}
\label{subsec:direction_data}
Die Rust-Simulation speichert für jeden Mikrofontreffer zusätzlich zur Ankunftszeit (\texttt{delay}) und Lautstärke (\texttt{pressure}) die Einfallsrichtung (\texttt{direction}) als 3D-Vektor.
Diese Richtungen werden in der parallelen Simulation gesammelt und in \texttt{ir\_output.json} exportiert.

\subsection{HRTF-Verarbeitung}
\label{subsec:hrtf_processing}
Für die binaurale Faltung wird eine HRTF-Datenbank (z.B. CIPIC SOFA-Datei) geladen.
Jede Reflexion wird anhand ihrer Richtung mit der passenden HRIR für linkes und rechtes Ohr gefaltet.
Das Ergebnis ist eine Stereo-WAV-Datei mit räumlichem Echo.

\subsection{Integration}
\label{subsec:hrtf_integration}
Die Python-Pipeline nutzt die Richtungsdaten aus \texttt{ir\_output.json}.
Falls Richtungen vorhanden sind, wird die HRTF-basierte Faltung verwendet, sonst die klassische.
